\section{Discretizzazione inversa e benchmark degli ottimizzatori}
\label{app:optimizer}
\setcounter{figure}{0}
\setcounter{table}{0}
\setcounter{equation}{0}

\subsection{Sensibilità del profilo di \texorpdfstring{$\beta$}{beta} alla griglia inversa}

Il target sintetico è stato mantenuto fisso a $(N,M)=(32,160)$, mentre la discretizzazione inversa è stata raffinata da $(16,80)$ a $(24,120)$. Il minimo campionato del profilo dati si sposta da $\beta=0.4500$ a $\beta=0.4375$. I parametri di disturbo $(\alpha,D,E)$ restano vicini tra le due risoluzioni; il termine dati al parametro vero diminuisce di un fattore $2.22$ e il minimo campionato del profilo di un fattore $3.53$ sulla griglia inversa più fine. Il confronto quantifica la sensibilità dei profili alla discretizzazione inversa. Le scansioni forward e backward sono procedure numeriche distinte; i valori riportati nel testo non vengono quindi interpretati come un’unica curva continua o come una prova di convergenza asintotica del problema inverso.


\subsection{Risultati per singola esecuzione}

Il benchmark comparativo impiega il medesimo target pulito congelato, la formulazione deterministica normalizzata (Modalità~A) con parametro di regolarizzazione $\lambda_A=10^{-4}$, centro a priori normalizzato $(0.5,0.5,0.5,0.5)$ e un budget massimo rigoroso di $600$ valutazioni effettive del solutore FEM per ciascuna esecuzione. Le tre esecuzioni CPSO provengono dalla campagna multi-seed precedente; hash del target, discretizzazioni e decomposizione del funzionale sono stati verificati rispetto ai relativi artefatti. Come dettagliato nella Tabella~\ref{tab:app-optimizer-runs}, una run di Differential Evolution si è arrestata a $588$ valutazioni per la tolleranza sulla dispersione della popolazione.

\begin{table}[H]
\centering
\caption{Risultati completi delle nove esecuzioni del benchmark a budget equivalente, con stime dei quattro parametri, costo totale finale e residuo osservativo relativo rispetto al target pulito. Il residuo è espresso in percentuale; la colonna delle valutazioni riporta il conteggio effettivo sul budget massimo pianificato.}
\label{tab:app-optimizer-runs}
\small
\setlength{\tabcolsep}{6.5pt}
\begin{tabular}{@{}lrrrrrrrr@{}}
\toprule
Metodo & Seed & Eval./600 & $\hat{\alpha}$ & $\hat{\beta}$ & $\hat{D}$ & $\hat{E}$ & $J_{\mathrm{total}}$ & $r_{\mathrm{clean}}$ [\%] \\
\midrule
CPSO                   & 20261017   & 600/600 & 0.7033 & 0.3880 & 0.07952 & 2.4788 & $6.3224\times10^{-6}$ & 0.2118 \\
CPSO                   & 20261018   & 600/600 & 0.6956 & 0.4791 & 0.07892 & 2.4805 & $4.4668\times10^{-6}$ & 0.1918 \\
CPSO                   & 20261019   & 600/600 & 0.6838 & 0.6033 & 0.07755 & 2.4870 & $8.9289\times10^{-6}$ & 0.5524 \\
\midrule
Differential Evolution & 2026120101 & 588/600 & 0.6964 & 0.4774 & 0.07901 & 2.4793 & $4.4626\times10^{-6}$ & 0.1851 \\
Differential Evolution & 2026120102 & 600/600 & 0.6964 & 0.4774 & 0.07901 & 2.4792 & $4.4626\times10^{-6}$ & 0.1851 \\
Differential Evolution & 2026120103 & 600/600 & 0.6964 & 0.4772 & 0.07901 & 2.4790 & $4.4627\times10^{-6}$ & 0.1848 \\
\midrule
Powell                 & 2026120201 & 600/600 & 0.6964 & 0.4774 & 0.07901 & 2.4794 & $4.4626\times10^{-6}$ & 0.1850 \\
Powell                 & 2026120202 & 600/600 & 0.6965 & 0.4773 & 0.07902 & 2.4791 & $4.4627\times10^{-6}$ & 0.1845 \\
Powell                 & 2026120203 & 600/600 & 0.6964 & 0.4774 & 0.07901 & 2.4793 & $4.4626\times10^{-6}$ & 0.1851 \\
\bottomrule
\end{tabular}
\end{table}

\subsection{Andamenti di convergenza e dispersione tra seed}

La storia di convergenza disponibile per la campagna multi-seed senza rumore è riportata nella Figura~\ref{fig:app-cpso-convergence-final}. Per i dieci twin experiments rumorosi sono disponibili riepiloghi finali e stime parametriche, ma non una cronologia completa verificabile del miglior candidato a ogni valutazione; non viene quindi rappresentata una traiettoria di convergenza per tali run. Il confronto diretto tra CPSO, Differential Evolution e il metodo di Powell basato su direzioni in termini di costo finale, residuo osservativo, recupero parametrico e traiettorie di discesa è illustrato nelle Figure~\ref{fig:app-optimizer-cost} e~\ref{fig:app-optimizer-convergence}.

La migliore esecuzione CPSO ($J_{\mathrm{total}}=4.4668\times10^{-6}$) raggiunge un costo vicino a quelli di Differential Evolution e Powell; la mediana CPSO è però più elevata e la dispersione tra inizializzazioni è maggiore, soprattutto per $\beta$. Il confronto descrive il comportamento dei metodi sul problema e sul budget considerati.

\begin{figure}[H]
\centering
\includegraphics[width=\linewidth]{03a_twin_noise_free_multistart/fig_multistart_convergence_v2.pdf}
\caption{Cronologia del miglior costo osservato nelle tre esecuzioni CPSO senza rumore. Tutte le run raggiungono il budget assegnato di 600 valutazioni; il grafico descrive la convergenza empirica di queste tre inizializzazioni e non certifica l’ottimalità globale.}
\label{fig:app-cpso-convergence-final}
\end{figure}

\begin{figure}[H]
\centering
\includegraphics[width=\linewidth]{12_optimizer_benchmark/fig_optimizer_benchmark_cost_and_errors.pdf}
\caption{Confronto descrittivo tra CPSO, Differential Evolution e Powell sul target sintetico congelato: costo finale e errori parametrici nel benchmark a tre repliche per metodo. Le differenze vanno interpretate nel quadro del budget e delle impostazioni riportati nel testo.}
\label{fig:app-optimizer-cost}
\end{figure}

\begin{figure}[H]
\centering
\includegraphics[width=\textwidth]{12_optimizer_benchmark/fig_optimizer_benchmark_convergence.pdf}
\caption{Traiettorie di convergenza delle nove esecuzioni del benchmark in funzione del numero di valutazioni effettive del solutore FEM. In questa campagna i valori finali raggiunti da Differential Evolution e Powell risultano più uniformi rispetto a CPSO; il risultato si riferisce al problema e alle impostazioni confrontate.}
\label{fig:app-optimizer-convergence}
\end{figure}